%=============================================================================
% CONTRARIAN ANALYSIS: The Case for kappa = 0 (Symmetric Split)
% Devil's Advocate Stress Test of the Constitutive Sector
% Generated: 2026-03-23
%=============================================================================

\documentclass[12pt,a4paper]{article}
\usepackage[margin=2.5cm]{geometry}
\usepackage{amsmath,amssymb,amsthm}
\usepackage{booktabs,array}
\usepackage{hyperref}
\usepackage{xcolor}
\usepackage{tcolorbox}

\newtheorem{theorem}{Theorem}
\newtheorem{proposition}{Proposition}
\newtheorem{claim}{Claim}
\newtheorem{objection}{Objection}
\theoremstyle{definition}
\newtheorem{definition}{Definition}
\newtheorem{response}{Response to Objection}
\theoremstyle{remark}
\newtheorem{remark}{Remark}

\newcommand{\CP}{\mathbb{C}P}
\newcommand{\dd}{\mathrm{d}}
\newcommand{\vE}{\mathbf{E}}
\newcommand{\vB}{\mathbf{B}}
\newcommand{\vD}{\mathbf{D}}
\newcommand{\vH}{\mathbf{H}}

\title{The Case for $\kappa = 0$:\\[6pt]
\large Why the Symmetric Constitutive Split May Be the Correct\\
DFD Prediction, and the Two Nonzero Values Are Artefacts}

\author{Contrarian Stress Test --- DFD Research Programme}
\date{23 March 2026}

\begin{document}
\maketitle

\begin{abstract}
The DFD vacuum loading paper introduces a constitutive split parameter
$\kappa$ governing the asymmetry between the electric permittivity
$\epsilon(\psi) = \epsilon_0\, n\, e^{+\kappa\psi}$ and magnetic
permeability $\mu(\psi) = \mu_0\, n\, e^{-\kappa\psi}$.  Two independent
derivation channels yield two \emph{different} nonzero values of $\kappa$,
creating an unresolved ``two-value problem.''  This document argues the
contrarian position: that $\kappa = 0$ (the symmetric split) is the
correct answer, both nonzero values are artefacts of additional assumptions
not forced by the minimal DFD postulates, and the symmetric-split theory
is simpler, stronger, and equally consistent with all observations.
The purpose of this document is adversarial stress-testing, not to
establish a conclusion.  Every argument here must be \emph{defeated}
before $\kappa \neq 0$ can be published with confidence.
\end{abstract}

\tableofcontents
\newpage

%=============================================================================
\section{The Setup: What $\kappa$ Controls}
\label{sec:setup}
%=============================================================================

In the DFD vacuum loading picture, the constitutive relations of
electrodynamics in a gravitational field are:
\begin{equation}
\epsilon(\psi) = \epsilon_0\, n(\psi)\, e^{+\kappa\psi}, \qquad
\mu_{\rm mag}(\psi) = \mu_0\, n(\psi)\, e^{-\kappa\psi}, \qquad
n = e^\psi.
\label{eq:constitutive}
\end{equation}
The phase velocity $v_{\rm ph} = 1/\sqrt{\epsilon\,\mu_{\rm mag}} = c/n$
is independent of $\kappa$ for \emph{all} values of $\kappa$.
Thus $\kappa$ controls only the \emph{relative} loading of the
electric versus magnetic sectors, not the speed of light.

Three values are in play:
\begin{center}
\renewcommand{\arraystretch}{1.3}
\begin{tabular}{lll}
\toprule
\textbf{Value} & \textbf{Source} & \textbf{Status} \\
\midrule
$\kappa = 0$ & E--B symmetry of the optical metric & Minimal assumption \\
$\kappa = \kappa_{\rm EW}$ & Electroweak derivation channel & Nonzero, specific value \\
$\kappa = \kappa_{\rm EH}$ & Einstein--Hilbert derivation channel & Nonzero, different value \\
\bottomrule
\end{tabular}
\end{center}

\begin{tcolorbox}[colback=red!5!white, colframe=red!60!black, title=Central Claim of This Document]
$\kappa = 0$ is the correct DFD prediction.  Both $\kappa_{\rm EW}$
and $\kappa_{\rm EH}$ involve auxiliary assumptions that go beyond the
minimal DFD postulates ($n = e^\psi$ plus the three axioms), and the
disagreement between them is evidence that these assumptions are
individually flawed.
\end{tcolorbox}

%=============================================================================
\section{Argument 1: The Optical Metric Symmetry}
\label{sec:optical-metric}
%=============================================================================

The DFD optical metric for photon propagation is:
\begin{equation}
\tilde{g}_{\mu\nu} = \mathrm{diag}\!\left(-\frac{c^2}{n^2},\; n^2,\; n^2,\; n^2\right).
\label{eq:optical-metric}
\end{equation}
This metric treats $\vE$ and $\vB$ on an identical footing.  Maxwell's
equations in this background are:
\begin{align}
\nabla \cdot (\epsilon_0\, n\, \vE) &= \rho_{\rm free}, &
\nabla \times \vE &= -\frac{\partial \vB}{\partial t}, \\
\nabla \cdot \vB &= 0, &
\nabla \times \left(\frac{\vB}{\mu_0\, n}\right) &= \vJ_{\rm free}
  + \frac{\partial(\epsilon_0\, n\, \vE)}{\partial t}.
\end{align}

\begin{claim}
Setting $\kappa = 0$ (i.e., $\epsilon = \epsilon_0\, n$ and
$\mu = \mu_0\, n$) is the \emph{unique} constitutive assignment that:
\begin{enumerate}
\item preserves the E--B duality symmetry of the free-field equations
  in the optical metric background;
\item has no preferred electromagnetic direction; and
\item reduces to the standard vacuum ($\epsilon_0, \mu_0$) when $\psi = 0$.
\end{enumerate}
\end{claim}

Any $\kappa \neq 0$ breaks the E--B symmetry of Eq.~\eqref{eq:optical-metric}.
Since the optical metric itself does not distinguish E from B, the
symmetry breaking must be \emph{introduced from outside}---it is not
a consequence of the DFD postulates alone.

%=============================================================================
\section{Argument 2: Maxwell Consistency at $\kappa = 0$}
\label{sec:maxwell-consistency}
%=============================================================================

Maxwell's equations in a medium with $\epsilon = \epsilon_0\, n$ and
$\mu = \mu_0\, n$ are perfectly consistent:
\begin{itemize}
\item The phase velocity is $v_{\rm ph} = c/n$ for \emph{both}
  polarizations (TE and TM).
\item There is no birefringence: the two polarization modes propagate
  at identical speeds.
\item The impedance is $Z = \sqrt{\mu/\epsilon} = \sqrt{\mu_0/\epsilon_0}
  = Z_0$, \emph{independent of $\psi$}.  This means there is no
  impedance mismatch between regions of different $\psi$---no
  reflection at the boundary of a gravitational well.
\item The constitutive relations are causal (Kramers--Kronig consistent)
  and satisfy the Schwarz reflection principle.
\end{itemize}

\begin{proposition}[No-go for birefringence detection at $\kappa = 0$]
If $\kappa = 0$, then to \emph{all orders} in $\psi$, there is no
gravitational birefringence, no polarization-dependent deflection, and
no difference between the propagation of TE and TM modes.
\end{proposition}

This is a \textbf{sharp, falsifiable prediction}: if any experiment
detects gravitational birefringence, $\kappa = 0$ is ruled out.
Conversely, every null detection of gravitational birefringence supports
$\kappa = 0$ against $\kappa \neq 0$.  Current observations are all
consistent with $\kappa = 0$.

%=============================================================================
\section{Argument 3: The Cassini--PPN Constraint}
\label{sec:cassini}
%=============================================================================

The Cassini spacecraft measurement constrains the PPN parameter $\gamma$:
\begin{equation}
|\gamma_{\rm PPN} - 1| < 2.3 \times 10^{-5}.
\end{equation}
DFD predicts $\gamma = 1$ exactly in the weak-field limit (confirmed
in the DFD corpus).

\begin{claim}
If $\kappa \neq 0$, the effective $\gamma$ parameter may differ for
electric-sector and magnetic-sector observables.  Specifically:
\begin{itemize}
\item The Shapiro time delay probes the photon propagation speed, which
  depends on $n$ but not on $\kappa$ (the phase velocity is $c/n$
  regardless of $\kappa$).
\item However, the gravitational deflection of light involves the
  \emph{gradient} of $n$, and if $\kappa \neq 0$, the effective
  refractive index experienced by different field components is
  $n_E = n\, e^{+\kappa\psi}$ versus $n_B = n\, e^{-\kappa\psi}$.
\item At first order in $\psi$, this produces a polarization-dependent
  correction $\delta\gamma \sim \kappa\psi \sim \kappa\,(GM/rc^2)$.
\end{itemize}
\end{claim}

\begin{remark}
For $\kappa \sim O(1)$ and solar-system $\psi \sim 10^{-6}$, the
correction $\delta\gamma \sim 10^{-6}$ is marginally within the Cassini
bound.  But for larger $\kappa$ or in stronger-field environments (neutron
star binaries, black hole shadows), the effect could become significant.
\textbf{Has a polarization-dependent PPN analysis been performed for DFD
with $\kappa \neq 0$?}  If not, $\kappa \neq 0$ carries an untested
observational liability.
\end{remark}

At $\kappa = 0$, there is no polarization dependence whatsoever, and
$\gamma = 1$ is guaranteed to be the same for all electromagnetic
observables.  This is the cleaner prediction.

%=============================================================================
\section{Argument 4: Tension with $c_T = c$}
\label{sec:tensor-speed}
%=============================================================================

DFD proves that gravitational waves propagate at exactly $c$:
\begin{equation}
c_T = c \quad \text{(exact)}.
\end{equation}
This is confirmed by GW170817 to $|c_T/c - 1| < 10^{-15}$.

The proof relies on the tensor perturbation $h_{ij}^{\rm TT}$ being
independent of $\psi$---the gravitational wave sector decouples from the
scalar sector.

\begin{claim}
If the EM sector has $\kappa \neq 0$, then the E and B components of
an EM wave in a gravitational background travel through media with
different effective indices.  Consider the limit where $\psi$ varies
slowly compared to the EM wavelength.  The electric field sees
$n_E = n\, e^{+\kappa\psi}$ and the magnetic field sees
$n_B = n\, e^{-\kappa\psi}$.  While the phase velocity of the combined
wave is $c/n$, the \emph{group velocities} of the E and B sectors
differ at $O(\kappa\, \nabla\psi)$.

This means that in a time-varying gravitational field (such as a passing
gravitational wave), the E and B components of a test EM wave respond
differently---their coupling to the GW background is
$\kappa$-dependent.  This introduces a polarization-dependent GW--EM
coupling that is \emph{absent} at $\kappa = 0$.
\end{claim}

\begin{remark}
The clean prediction $c_T = c$ is one of DFD's crown jewels.  Introducing
$\kappa \neq 0$ potentially complicates the EM sector's response to
gravitational waves, even if $c_T$ itself (the speed of the GW) is
unaffected.  The question is whether the EM response to GW is also
``clean.''  At $\kappa = 0$ it is, automatically.
\end{remark}

%=============================================================================
\section{Argument 5: $\kappa \neq 0$ Requires Additional Physics}
\label{sec:additional-physics}
%=============================================================================

The minimal DFD postulate set is:
\begin{enumerate}
\item A real scalar field $\psi$ on spacetime.
\item The refractive index law $n = e^\psi$ (from Cauchy's functional
  equation and the multiplicative composition requirement).
\item The three DFD axioms (gauge emergence, topological truncation,
  and the master invariant).
\end{enumerate}

\begin{claim}
$\kappa \neq 0$ is \textbf{not a consequence} of these postulates.
\end{claim}

Both derivation channels that produce $\kappa \neq 0$ introduce
additional structure:
\begin{itemize}
\item \textbf{Electroweak channel:} Assumes a specific coupling mechanism
  between the $\psi$-field and the electroweak gauge sector that
  differentiates between the electric and magnetic response.  This
  coupling is \emph{not} a minimal consequence of $n = e^\psi$; it
  requires specifying \emph{how} the gauge fields source $\psi$ at the
  level of individual field components.
\item \textbf{Einstein--Hilbert channel:} Assumes a specific form of the
  gravitational action that distinguishes the E and B sectors when
  expanded to second order in $\psi$.  This is a statement about the
  \emph{gravitational back-reaction} of EM fields, not about the
  refractive index law.
\end{itemize}

The fact that these two channels produce \emph{different} values of
$\kappa$ is strong evidence that the additional assumptions in each
channel are not uniquely determined by the DFD axioms.  If $\kappa$
were a fundamental DFD prediction, both channels should agree.

\begin{tcolorbox}[colback=yellow!5!white, colframe=yellow!60!black, title=The Two-Value Problem as Evidence]
The two-value problem is not merely an unresolved technical issue.
It is \emph{diagnostic evidence} that $\kappa$ is an artefact of
assumptions beyond the minimal DFD framework.  The resolution is not
to ``fix'' one channel or the other, but to recognise that $\kappa = 0$
requires \emph{no} additional assumptions and is the unique value
consistent with the minimal postulates.
\end{tcolorbox}

%=============================================================================
\section{Argument 6: Occam's Razor}
\label{sec:occam}
%=============================================================================

\begin{center}
\renewcommand{\arraystretch}{1.3}
\begin{tabular}{lcc}
\toprule
\textbf{Property} & $\kappa = 0$ & $\kappa \neq 0$ \\
\midrule
Free parameters in constitutive sector & 0 & 1 \\
E--B symmetry preserved & Yes & No \\
Birefringence prediction & Null (clean) & Nonzero (complex) \\
Impedance $Z(\psi)$ & $Z_0$ (constant) & $Z_0\, e^{-\kappa\psi}$ \\
Two-value problem & Absent & Present \\
Additional physics beyond $n = e^\psi$ & None required & Required \\
Consistent with all current observations & Yes & Yes \\
Attack surface for critics & Minimal & Significant \\
\bottomrule
\end{tabular}
\end{center}

By every metric of parsimony and theoretical hygiene, $\kappa = 0$ is
the superior choice.  It introduces no unexplained parameter, requires
no additional derivation to justify, creates no two-value puzzle, and
matches all existing data.

%=============================================================================
\section{Addressing Counterarguments}
\label{sec:counterarguments}
%=============================================================================

\subsection{``Vacuum Loading Requires $\kappa \neq 0$ Because EM Energy
  Loads the Vacuum''}

\begin{objection}
The physical picture is that electromagnetic energy is stored in the
vacuum, and the electric and magnetic sectors store energy differently
because $\vE$ and $\vB$ couple to different geometric structures.
This forces $\kappa \neq 0$.
\end{objection}

\begin{response}
Vacuum loading can be entirely symmetric.  In the standard
electromagnetic energy density:
\begin{equation}
u_{\rm EM} = \frac{1}{2}\left(\epsilon\, |\vE|^2 + \frac{|\vB|^2}{\mu}\right),
\end{equation}
setting $\epsilon = \epsilon_0\, n$ and $\mu = \mu_0\, n$ gives:
\begin{equation}
u_{\rm EM} = \frac{n}{2}\left(\epsilon_0\, |\vE|^2
  + \frac{|\vB|^2}{\mu_0}\right)
  = n \cdot u_{\rm EM}^{(\rm flat)}.
\end{equation}
The loading factor is $n$ for \emph{both} sectors.  There is no
asymmetry.  The claim that loading must be asymmetric is an \emph{additional
assumption}, not a consequence of the loading picture.

Moreover, this symmetric loading has a transparent physical interpretation:
the vacuum is a single medium with a single refractive index $n$, and
both $\epsilon$ and $\mu$ are modified by the same factor $n$.  This is
precisely what happens in a non-dispersive, non-birefringent dielectric.
The vacuum loading picture, at its most natural, \emph{predicts}
$\kappa = 0$.
\end{response}

\subsection{``The Paper Predicts Testable Consequences of $\kappa \neq 0$''}

\begin{objection}
The vacuum loading paper derives specific experimental signatures of
$\kappa \neq 0$, including modifications to Coulomb's law, the LPI
observable $\xi(\kappa)$, and gravitational birefringence.  These
predictions give the theory teeth.
\end{objection}

\begin{response}
The predictions for $\kappa = 0$ are \emph{equally testable}---they
are null predictions:
\begin{itemize}
\item No modification to Coulomb's law beyond the overall $n$-scaling.
\item $\xi = 0$ (no E--B asymmetry in local position invariance tests).
\item No gravitational birefringence at any level.
\item Constant impedance $Z = Z_0$ everywhere.
\end{itemize}
A null prediction is not a non-prediction.  It is a sharp, specific,
falsifiable claim.  Every experiment that tests for $\kappa$-dependent
effects and finds null is a \emph{confirmation} of $\kappa = 0$.

Furthermore, $\kappa = 0$ predictions are \textbf{more falsifiable}
in a Popperian sense: a single detection of gravitational birefringence
would definitively rule out $\kappa = 0$, whereas $\kappa \neq 0$
can always adjust the value of $\kappa$ to accommodate a non-detection
(``$\kappa$ is small enough to evade the bound'').
\end{response}

\subsection{``But the Electroweak Derivation Is Rigorous''}

\begin{objection}
The electroweak channel derives $\kappa$ from the gauge--$\psi$ coupling
in a rigorous way.  You cannot simply ignore it.
\end{objection}

\begin{response}
If the electroweak derivation were the unique, unavoidable consequence
of the DFD axioms, it would have to agree with the Einstein--Hilbert
channel.  It does not.  This means at least one of the following is true:
\begin{enumerate}
\item The electroweak channel makes an assumption not shared by the
  EH channel.
\item The EH channel makes an assumption not shared by the electroweak
  channel.
\item Both channels make different auxiliary assumptions.
\end{enumerate}
In all three cases, $\kappa \neq 0$ is \emph{not uniquely determined}
by the DFD axioms.  The disagreement is proof that additional input is
required.  $\kappa = 0$, by contrast, requires \emph{no} auxiliary input.
\end{response}

%=============================================================================
\section{What $\kappa = 0$ Does for the Vacuum Loading Paper}
\label{sec:simplification}
%=============================================================================

If $\kappa = 0$ is adopted:
\begin{enumerate}
\item The entire ``constitutive split'' section can be \textbf{removed}
  or replaced by a one-paragraph statement that E--B symmetry is
  preserved.
\item The two-value problem \textbf{disappears}.
\item The parameter count drops by one.
\item The observational predictions become simpler and cleaner.
\item The paper's attack surface shrinks: critics cannot point to the
  two-value disagreement as evidence of internal inconsistency.
\item The connection to the optical metric~\eqref{eq:optical-metric}
  becomes tighter: the constitutive relations are the \emph{unique}
  choice consistent with the optical metric's E--B symmetry.
\end{enumerate}

The paper becomes \textbf{shorter}, \textbf{simpler}, and
\textbf{harder to attack}.

%=============================================================================
\section{What Must Be True for $\kappa \neq 0$ to Survive}
\label{sec:survival-conditions}
%=============================================================================

To defeat this contrarian argument and establish $\kappa \neq 0$, the
following would all need to be demonstrated:

\begin{enumerate}
\item \textbf{Resolution of the two-value problem:} Either show that
  both channels give the same $\kappa$ (perhaps the claimed disagreement
  is a computational error) or show rigorously that one channel is
  wrong and why.

\item \textbf{Derivation from minimal postulates:} Show that $\kappa = 0$
  is actually \emph{inconsistent} with the DFD axioms, or that
  $\kappa \neq 0$ follows inevitably from $n = e^\psi$ without
  additional assumptions.

\item \textbf{Polarization-dependent PPN analysis:} Perform the
  full PPN analysis with $\kappa \neq 0$ and show that the
  polarization-dependent corrections to $\gamma$ are consistent with
  all solar-system tests (Cassini, lunar laser ranging, VLBI).

\item \textbf{EM--GW coupling analysis:} Show that $\kappa \neq 0$ does
  not introduce observable polarization-dependent effects in the
  response of electromagnetic waves to gravitational waves.

\item \textbf{Physical motivation:} Provide a clear, minimal physical
  argument for \emph{why} the vacuum should load E and B differently,
  beyond the fact that two formal channels happen to produce $\kappa \neq 0$.
\end{enumerate}

Until all five conditions are met, $\kappa = 0$ remains the
parsimonious, safe, and defensible choice.

%=============================================================================
\section{Summary}
\label{sec:summary}
%=============================================================================

\begin{tcolorbox}[colback=blue!5!white, colframe=blue!60!black, title=Bottom Line]
The symmetric constitutive split $\kappa = 0$ is:
\begin{itemize}
\item Forced by E--B symmetry of the DFD optical metric.
\item Consistent with Maxwell's equations without birefringence.
\item Compatible with all current observations (Cassini, GW170817,
  null birefringence searches).
\item The unique value requiring no physics beyond the minimal DFD
  postulates.
\item Free of the two-value problem that plagues $\kappa \neq 0$.
\item Simpler, cleaner, and harder to attack.
\end{itemize}

The two nonzero values of $\kappa$ are artefacts of auxiliary
assumptions introduced by the electroweak and Einstein--Hilbert
derivation channels, respectively.  Their mutual disagreement is
evidence that these assumptions are not consequences of the DFD axioms.

\textbf{Recommendation:} Either resolve all five survival conditions
in Section~\ref{sec:survival-conditions}, or adopt $\kappa = 0$ and
simplify the vacuum loading paper accordingly.
\end{tcolorbox}

\end{document}
